% fig06_training_pipeline.tex — 图6 LLM 训练流水线三阶段
% 《深度学习与大语言模型：前沿技术全景报告》图示库
% 由 dl_llm_report.tex 抽取，经 % fig06_training_pipeline.tex — 图6 LLM 训练流水线三阶段
% 《深度学习与大语言模型：前沿技术全景报告》图示库
% 由 dl_llm_report.tex 抽取，经 % fig06_training_pipeline.tex — 图6 LLM 训练流水线三阶段
% 《深度学习与大语言模型：前沿技术全景报告》图示库
% 由 dl_llm_report.tex 抽取，经 % fig06_training_pipeline.tex — 图6 LLM 训练流水线三阶段
% 《深度学习与大语言模型：前沿技术全景报告》图示库
% 由 dl_llm_report.tex 抽取，经 \input{figs/fig06_training_pipeline} 引用（caption/label 在主文件）
% 注意：本文件为片段，依赖主文件导言区的 tikz/pgfplots 宏包、颜色定义与 tikzset 样式，不可单独编译。

\begin{tikzpicture}[node distance=9mm]
\node[blk, minimum width=46mm, minimum height=23mm, align=center] (pre) {\textbf{阶段一：预训练}\\ 万亿级 token / 万卡集群\\ 目标：下一 token 预测\\ {\scriptsize 成本：百万 GPU$\cdot$小时级}};
\node[blk, right=10mm of pre, minimum width=42mm, minimum height=23mm, align=center] (sft) {\textbf{阶段二：监督微调}\\ $10^4$--$10^6$ 条指令-回答对\\ 目标：遵循指令\\ {\scriptsize 成本：数千 GPU$\cdot$小时级}};
\node[blk, right=10mm of sft, minimum width=42mm, minimum height=23mm, align=center] (rl) {\textbf{阶段三：对齐}\\ RLHF / DPO / RLVR\\ 目标：有用且无害\\ {\scriptsize 成本：数万 GPU$\cdot$小时级起}};
\node[gry, below=9mm of sft, minimum width=104mm, minimum height=8mm] (cap) {能力演进：语言与知识 $\to$ 指令遵循 $\to$ 偏好对齐与推理};
\draw[arr] (pre)--(sft); \draw[arr] (sft)--(rl);
\draw[arr] (rl.south) -- ++(0,-20mm) -| node[lab, pos=0.25, below] {数据飞轮：用户交互 $\to$ 高质量数据回流} (pre.south);
\end{tikzpicture}
 引用（caption/label 在主文件）
% 注意：本文件为片段，依赖主文件导言区的 tikz/pgfplots 宏包、颜色定义与 tikzset 样式，不可单独编译。

\begin{tikzpicture}[node distance=9mm]
\node[blk, minimum width=46mm, minimum height=23mm, align=center] (pre) {\textbf{阶段一：预训练}\\ 万亿级 token / 万卡集群\\ 目标：下一 token 预测\\ {\scriptsize 成本：百万 GPU$\cdot$小时级}};
\node[blk, right=10mm of pre, minimum width=42mm, minimum height=23mm, align=center] (sft) {\textbf{阶段二：监督微调}\\ $10^4$--$10^6$ 条指令-回答对\\ 目标：遵循指令\\ {\scriptsize 成本：数千 GPU$\cdot$小时级}};
\node[blk, right=10mm of sft, minimum width=42mm, minimum height=23mm, align=center] (rl) {\textbf{阶段三：对齐}\\ RLHF / DPO / RLVR\\ 目标：有用且无害\\ {\scriptsize 成本：数万 GPU$\cdot$小时级起}};
\node[gry, below=9mm of sft, minimum width=104mm, minimum height=8mm] (cap) {能力演进：语言与知识 $\to$ 指令遵循 $\to$ 偏好对齐与推理};
\draw[arr] (pre)--(sft); \draw[arr] (sft)--(rl);
\draw[arr] (rl.south) -- ++(0,-20mm) -| node[lab, pos=0.25, below] {数据飞轮：用户交互 $\to$ 高质量数据回流} (pre.south);
\end{tikzpicture}
 引用（caption/label 在主文件）
% 注意：本文件为片段，依赖主文件导言区的 tikz/pgfplots 宏包、颜色定义与 tikzset 样式，不可单独编译。

\begin{tikzpicture}[node distance=9mm]
\node[blk, minimum width=46mm, minimum height=23mm, align=center] (pre) {\textbf{阶段一：预训练}\\ 万亿级 token / 万卡集群\\ 目标：下一 token 预测\\ {\scriptsize 成本：百万 GPU$\cdot$小时级}};
\node[blk, right=10mm of pre, minimum width=42mm, minimum height=23mm, align=center] (sft) {\textbf{阶段二：监督微调}\\ $10^4$--$10^6$ 条指令-回答对\\ 目标：遵循指令\\ {\scriptsize 成本：数千 GPU$\cdot$小时级}};
\node[blk, right=10mm of sft, minimum width=42mm, minimum height=23mm, align=center] (rl) {\textbf{阶段三：对齐}\\ RLHF / DPO / RLVR\\ 目标：有用且无害\\ {\scriptsize 成本：数万 GPU$\cdot$小时级起}};
\node[gry, below=9mm of sft, minimum width=104mm, minimum height=8mm] (cap) {能力演进：语言与知识 $\to$ 指令遵循 $\to$ 偏好对齐与推理};
\draw[arr] (pre)--(sft); \draw[arr] (sft)--(rl);
\draw[arr] (rl.south) -- ++(0,-20mm) -| node[lab, pos=0.25, below] {数据飞轮：用户交互 $\to$ 高质量数据回流} (pre.south);
\end{tikzpicture}
 引用（caption/label 在主文件）
% 注意：本文件为片段，依赖主文件导言区的 tikz/pgfplots 宏包、颜色定义与 tikzset 样式，不可单独编译。

\begin{tikzpicture}[node distance=9mm]
\node[blk, minimum width=46mm, minimum height=23mm, align=center] (pre) {\textbf{阶段一：预训练}\\ 万亿级 token / 万卡集群\\ 目标：下一 token 预测\\ {\scriptsize 成本：百万 GPU$\cdot$小时级}};
\node[blk, right=10mm of pre, minimum width=42mm, minimum height=23mm, align=center] (sft) {\textbf{阶段二：监督微调}\\ $10^4$--$10^6$ 条指令-回答对\\ 目标：遵循指令\\ {\scriptsize 成本：数千 GPU$\cdot$小时级}};
\node[blk, right=10mm of sft, minimum width=42mm, minimum height=23mm, align=center] (rl) {\textbf{阶段三：对齐}\\ RLHF / DPO / RLVR\\ 目标：有用且无害\\ {\scriptsize 成本：数万 GPU$\cdot$小时级起}};
\node[gry, below=9mm of sft, minimum width=104mm, minimum height=8mm] (cap) {能力演进：语言与知识 $\to$ 指令遵循 $\to$ 偏好对齐与推理};
\draw[arr] (pre)--(sft); \draw[arr] (sft)--(rl);
\draw[arr] (rl.south) -- ++(0,-20mm) -| node[lab, pos=0.25, below] {数据飞轮：用户交互 $\to$ 高质量数据回流} (pre.south);
\end{tikzpicture}
